\SourcePage{276}{281}
\section{Resonance fluorescence}

Taking into account the finite widths of levels in the problem of scattering of light is essential in cases when the frequency $\omega$ of the incident light is close to one of the ``intermediate'' frequencies $\omega_{n1}$ or $\omega_{n2}$---the so-called \textit{resonance fluorescence} (\textit{V.\,Weisskopf}, 1931).

Let us consider unshifted scattering by a system (say, an atom) in the ground state, so that the initial and final levels coincide and are strictly discrete. Let the frequency of the light be close to a certain frequency $\omega_{n1}$, where the level~$n$ is excited and therefore quasidiscrete.

This question could be solved by the method presented in the preceding paragraph. In this, however, there is no necessity, since the problem is completely analogous to that of the non-relativistic resonance scattering considered in vol.\,III, \S\,134. According to the results obtained there the resonance amplitude of scattering must contain the pole factor
\[
\left[\omega - \left(E_n - i\,\frac{\Gamma_n}{2}\right) - E_1\right]^{-1}.
\]

On the other hand, for $|\omega - \omega_{n1}| \gg \Gamma_n$ the formula must go over to the non-resonance formula~(59.5). From this it is clear that the desired cross section of scattering is obtained simply by the replacement of $E_n$ by $E_n - i\Gamma_n/2$ in formula~(59.5), and in the sum over~$n$ one may limit oneself to the resonance terms only:
\begin{equation}
d\sigma = \frac{\left|\sum_{M_n}(\mathbf{d}_{2n}\mathbf{e}'^*)(\mathbf{d}_{n1}\mathbf{e})\right|^2}{(\omega_{n1} - \omega)^2 + \Gamma^2_n/4}\;\omega^4\,do'.
\tag{63.1}
\end{equation}

The summation is carried out over all states (with different projections $M_n$), corresponding to the resonance level~$E_n$; in the sum over~$n$ in formula~(59.5) the states 1 and~2 may differ by different values of $M_1$ and $M_2$.

The cross section~(63.1) is maximal at $\omega = \omega_{n1}$. By order of magnitude the value at the maximum is $\sigma_\mathrm{max} \sim \omega^4d^4/\Gamma^2_n$. Since the probability of spontaneous transition $n \to 1$, and with it also the width~$\Gamma_n$, are of order $\omega^3d^2$, this value
\begin{equation}
\sigma_\mathrm{max} \sim \omega^{-2} \sim \lambda^2,
\tag{63.2}
\end{equation}
i.e.\ of order of the square of the wavelength of the light and does not depend on the constant of the fine structure---in place of the typical values $\sim r^2_e$ outside the region of resonance.

Let us emphasise that since the atom both before and after scattering is on a strictly discrete (ground) level, the frequencies of the primary and secondary photons strictly coincide. Therefore in

\SourcePage{277}{282}
irradiation with monochromatic light the scattered light will also be monochromatic. If, however, the incident light has a spectral distribution of intensities $I(\omega)$, and the function $I(\omega)$ changes little over the width~$\Gamma_n$, then the intensity of the scattered light will be proportional to
\[
\frac{I(\omega_{n1})\,d\omega}{(\omega - \omega_{n1})^2 + \Gamma^2_n/4}.
\tag{63.3}
\]
In other words, the line shape of the scattered light will coincide with the natural line shape in spontaneous emission from the level~$E_n$.

The cross section~(63.1) corresponds to the scattering tensor
\begin{equation}
(c_{ik})_{21} = \frac{\sum_{M_n}(d_i)_{2n}(d_k)_{n1}}{\omega_{n1} - \omega - i\Gamma_n/2}.
\tag{63.4}
\end{equation}
In particular, the tensor of polarisability
\begin{equation}
\frac{\sum_{M_n}(d_i)_{1n}(d_k)_{n1}}{\omega_{n1} - \omega - i\Gamma_n/2}.
\tag{63.5}
\end{equation}

Let us at once note that the addition of the imaginary part to the energies of intermediate excited states violates the Hermiticity of the tensor of polarisability and also at frequencies below the ionisation threshold. Its imaginary part appears, directly associated with the absorption of light. Having absorbed a quantum, the atom will sooner or later return to the ground state with emission of one or several photons. Therefore from the point of view of the absorption cross section the absorption cross section is simply the total cross section $\sigma_t$ of all possible scattering processes\,\footnote{Let us emphasise that we are speaking of the absorption by a system in a stable, ground state. On account of the finiteness of the lifetime of the excited state the formulation of the question would be different.}. On the other hand, according to formula~(59.25), expressing the optical theorem, this cross section is determined by the anti-Hermitian part of the tensor of polarisability.

Substituting in~(59.25) the tensor $\alpha_{ik}$ from~(63.5), we find the following formula for the absorption cross section of a photon of frequency $\omega$ close to~$\omega_{n1}$:
\begin{equation}
\sigma^{(\mathrm{abs})} = 4\pi^2\sum_{M_n}|\mathbf{d}_{n1}\mathbf{e}|^2\omega\,\frac{\Gamma_n/2}{\pi[(\omega - \omega_{n1})^2 + \Gamma^2_n/4]}.
\tag{63.6}
\end{equation}

\SourcePage{278}{283}
In the limit $\Gamma_n \to 0$ the last factor in this formula tends to the $\delta$-function $\delta(\omega - \omega_{n1})$, in accordance with the fact that in this case only a photon of a strictly defined frequency can be absorbed. Let light fall on the atom with spectral and angular distribution given by the flux density of energy $I_{\mathbf{ke}}$ (cf.~(44.7)). Then the photon number flux density is $I_{\mathbf{ke}}/\omega$, and the absorption probability
\begin{equation}
dw^{(\mathrm{abs})} = \sigma^{(\mathrm{abs})}\,\frac{I_{\mathbf{ke}}}{\omega}\,d\omega\,do.
\tag{63.7}
\end{equation}
If the function $I_{\mathbf{ke}}(\omega)$ changes little over the width~$\Gamma_n$, then after integrating over frequencies we obtain
\[
dw^{(\mathrm{abs})} = 4\pi^2\sum_{M_n}|\mathbf{d}_{n1}\mathbf{e}|^2\,I_{\mathbf{ke}}(\omega_{n1})\,d\omega\,do.
\]

Noting, on the other hand, that according to~(45.5),
\[
dw^{(\mathrm{sp})} = \frac{\omega^3}{2\pi}\sum_{M_n}|\mathbf{d}_{1n}\mathbf{e}^*|^2\,do = \frac{\omega^3}{2\pi}\sum_{M_n}|\mathbf{d}_{n1}\mathbf{e}|^2\,do
\]
is the probability of spontaneous emission of a photon of frequency~$\omega_{n1}$, we return to formula~(44.9).
